\section{Integral elimination by a finite B\'ezout identity}
\label{sec:integral-elimination}

There is a direct construction of the row operation for \(J_n\).
Its coefficients use a single univariate B\'ezout identity for each height.
Write \(z=XY\), \(N=k+1\), and retain \(B_2(z)=z\) and the
factors \(B_d(z)\) defined above. For \(2\le r\le N\) put
\[
 t_{N,r}(z)=\prod_{\substack{2\le d\le N\\d\nmid r}}B_d(z),
 \qquad
 I_N=\{r:2\le r\le N,\ 2r>N\}.
\]

\begin{lemma}
There exist \(a_{N,r}(z)\in\mathbb Z[z]\) such that
\(\sum_{r\in I_N}a_{N,r}(z)t_{N,r}(z)=1\).
\end{lemma}
\begin{proof}
For incomparable indices \(d,e\) under divisibility the ideal
\((B_d,B_e)\) is the unit ideal of \(\mathbb Z[z]\).
Indeed specialize \(x_s=z,x_t=1\) in Hazi's
two-colour B\'ezout identity for the primitive factors of colour \(s\)
\cite[Lemma 3.5]{hazi}. This also covers an index equal to two.
Multiply one such identity over every incomparable pair from
\(\{2,\ldots,N\}\). In each expanded summand the indices whose
factors were never selected form a chain. Choose \(r\in I_N\) divisible
by the largest member of the chain, or any \(r\in I_N\) if the chain
is empty. Such a multiple exists by choosing the largest multiple
up to \(N\). Each \(d\nmid r\) has therefore been selected at least
once, so the summand is a multiple of \(t_{N,r}\). Group the summands.
The same reasoning applies to the empty product of identities.
\end{proof}

\begin{lemma}
For a row of height \(k\) and each \(r\in I_N\), one can add an
integral polynomial combination of rows of strictly smaller area
to \(t_{N,r}(XY)v_P\) to obtain an element of
\(b_k(X,Y)\mathbb Z[X,Y]^{|\mathcal P_n|}\).
\end{lemma}
\begin{proof}
Work first in the generic two-colour diagram calculus. For either
orientation \(u\), a common denominator of the coefficients of
the Jones--Wenzl projector \(f^u_{r-1}\) is
\[
 g_{r-1,u}=\prod_{\substack{2\le d<r\\d\nmid r}}\Theta_{d,u}.
\]
This is the denominator statement in Hazi's proof of Theorem A
\cite[\S4, Lemma 4.2 and the proof of Theorem A]{hazi}.
Under \(x_s=-X, x_t=-Y\) the factors with \(d>2\) give
\(B_d(XY)\). An index \(d=2\) contributes either \(-X\) or \(-Y\),
both of which divide \(B_2(XY)\). Hence \(g_{r-1,u}\mid t_{N,r}(XY)\).
The box \(t_{N,r}(XY)f^u_{r-1}\) is an integral polynomial
linear combination of diagrams, with identity coefficient \(t_{N,r}(XY)\).

Insert this box on \(r-1\) consecutive arcs crossing a maximal gap in
\(P\). Every nonidentity term reduces the number crossing that gap.
At any earlier gap, only the selected endpoints already encountered
can contribute a crossing of a redrawn arc, and each contributed one
previously. Use the reversed order at later gaps. Thus the remaining
terms have smaller area. Each selected half-arc leads to a distinct
marked boundary endpoint, so expanding the box creates no closed
circle on the row side alone.

Fix the leftmost gap of maximum height in each row \(P\). At that
gap order the crossing arcs from outermost to innermost. Use the
first \(m=r-1\) for the insertion, so that the choices for different
\(r\) are nested. Write \(K_{P,Q,r}\) for the value of the entry
after inserting \(t_{N,r}(XY)f^u_m\) and closing with \(Q\).
Put \(z=XY\) and
\[
                 E_r(z)=\prod_{\substack{d\mid r\\ d>2}} B_d(z).
\]

The complement of the box in the sphere is a disk. It contains
circles and a matching of the \(2m\) ports. Delete the circles starting
with any innermost disk on the side away from the box. Such a disk
contains no port or arc ending at a port. The deletion contributes
the colour weight of its interior face. The resulting monomial \(C\)
is common to all terms of the projector and to the identity box:
each term has the same prescribed boundary shading. After all these
deletions, an arc returning to the same row of ports has an innermost
adjacent turnback on that row, so the value with the projector is
zero. This uses the turnback relation over the fraction field;
the cleared box already has coefficients in the polynomial ring.

If every complementary arc passes from one row of ports to the other,
noncrossing forces the unique order-preserving matching. Choose
the exterior point next to the left edge of the box and draw this
matching as a right closure. Let \(w\in\{X,Y\}\) be the weight of the
left face, and \(w'\) the opposite weight, so that \(ww'=z\).
For the left boundary colour \(u\), Hazi's disk rule gives
\([2]_u=-w'\). Closing the rightmost strand keeps the left boundary
colour fixed. Iterating the partial trace identity
\(\operatorname{pTr}(f_j^u)=-[j+1]_u f_{j-1}^u/[j]_u\)
in the fraction field gives the planar scalar \((-1)^m[r]_u\)
\cite[equation (7)]{hazi}. Multiplying by \(w\) counts the outer
face once. The spherical values of the boxes alone are therefore
\[
 W_f=
 \begin{cases}
       z E_r(z),&r\ \text{even},\\
       w E_r(z),&r\ \text{odd},
 \end{cases}
 \qquad
 W_{\rm id}=w^{\lfloor m/2\rfloor+1}(w')^{\lceil m/2\rceil}.
\]
Indeed \([r]_u=E_r(z)\) for odd \(r\) and
\([r]_u=-w' E_r(z)\) for even \(r\). For the identity, the \(m\)
nested circles give alternating face colours starting with \(w\).
Choosing the exterior point on the other side gives the same two
spherical values: the side colours coincide for even \(m\), and both
expressions are symmetric in them for odd \(m\). The boundary ports
of the box stay fixed throughout. With the outside of the stack on
the left, \(w=Y\) for the chosen prefix: the outer face of the
matching meets the white interval from \(2n-1\) to \(0\).

Write \(a=|\pi\vee\rho|\) and \(a'=|\pi'\vee\rho'|\). Restoring
the identity box recovers the original pair, hence
\(C W_{\rm id}=J_n(P,Q)=X^aY^{a'}\). Combining the two values with
the definitions of \(t_{N,r}\) and \(b_k\) gives the exact quotient
\begin{equation}
\frac{K_{P,Q,r}}{b_k(X,Y)}
=
\begin{cases}
 X^{a-\lfloor r/2\rfloor}Y^{a'-\lfloor r/2\rfloor},
   &\text{the selected arcs lie in distinct cycles of }P\cup Q,\\
 0,&\text{otherwise}.
\end{cases} \label{eq:through-quotient}
\end{equation}
To see the criterion, cut the \(m\) selected arcs in the identity
diagram. When they belong to distinct cycles, the two ends of each
one are connected by the remainder of its own cycle. Conversely, the
unique through matching reconnects precisely those pairs, so
restoring the identity puts each selected arc on its own cycle.
In the through case the displayed quotient equals \(C\) for even
\(r\) and \(wC\) for odd \(r\); its exponents are nonnegative.
This proves the required integral divisibility for every pair.

\end{proof}

Multiply the rows of the second lemma by the coefficients of the
first lemma and add. The coefficient of \(v_P\) is exactly one.
Doing this for every \(P\), in increasing order of area, constructs
a lower unitriangular \(L\) with \(LJ_n=DV\), where \(V\) has
integral polynomial entries. The degree count in Section~\ref{sec:degree}
gives \(\det V=1\). Multiplying by \(V^{-1}\) gives the claimed
diagonal; simultaneous signed permutations sort the heights.
For \(k=3\), the scalar identity in this construction is simply
\[
             1=(z-1)^2-z(z-2),\qquad I_4=\{3,4\}.
\]
Finally, at \(X=Y=x\) the spherical region count is the loop count
plus one. Thus \(J_n(x,x)=xG_n(x)\), while
\(b_k(x,x)=x\,s_k(x)\). Cancelling \(x\) in the matrix identity
gives the one-variable special Smith form, preserving both
determinants. Explicitly, its entries are \(s_{h(P)}(x)\) in
nondecreasing order of height. This is Conjecture~2. The implication
from the two-variable form was already recorded by Miller and
Stanton \cite[arXiv version, equation (64)]{ms}.

\begin{corollary}\label{cor:entries}
Fix the row order, the outermost prefixes and the B\'ezout weights above.
The entries of the right multiplier in \(LJ_n=DV\) are determined
by four integers. For \(P,Q\) write \(k=h(P)\), and let \(q\) be
the length of the longest initial segment of the \(k\) chosen
crossing arcs that belong to pairwise distinct cycles of \(P\cup Q\).
Let \(a,a'\) be the two block counts in \(J_n(P,Q)\) and set
\(\ell=\lfloor(q+1)/2\rfloor\). Then
\[
 V_{P,Q}=X^{a-\ell}Y^{a'-\ell}\,
 \sum_{\substack{r\in I_{k+1}\\r\le q+1}}
   a_{k+1,r}(XY)(XY)^{\,\ell-\lfloor r/2\rfloor}.
\]
\end{corollary}
\begin{proof}
In the weighted sum of \eqref{eq:through-quotient} the surviving
prefixes have length at most \(q\). Applying the identity-box face
count to the first \(q\) arcs gives \(a,a'\ge\ell\), whether or not
\(q+1\) belongs to \(I_{k+1}\). Factoring out the displayed
monomial gives the result.
\end{proof}
Thus at order \(n\) one can precompute at most \(n(n+1)/2\) univariate
polynomials, one for each pair \(1\le q\le k\le n\).
For a particular pair of matchings, a traversal of its alternating
cycles determines \(q\). If the total number of cycles is less than
\(\lfloor(k+1)/2\rfloor\), the sum is empty. In particular a connected
pair gives a zero entry of \(V\) whenever the row height is at least
three. These statements concern the change of basis.

